% !TEX program = xelatex
\documentclass[11pt,a4paper]{ctexart}
\usepackage[margin=1in]{geometry}
\usepackage[hidelinks]{hyperref}
\usepackage{multicol}
\usepackage{titlesec}

% Section formatting
\titleformat*{\section}{\large\bfseries}
\titleformat*{\subsection}{\normalsize\bfseries}

\begin{document}

%==== Header ====
\begin{center}
  {\LARGE \textbf{苏彦颀}} \\
  慕尼黑工业大学 - 信息与计算技术学院 \\
  \href{mailto:yanqi.su@tum.de}{yanqi.su@tum.de} \quad (+49) 151 58748540
\end{center}
\vspace{0.5em}
\hrule
\vspace{1em}

%==== Sections ====

\section*{研究兴趣}
软件测试；软件工程知识图谱；软件仓库挖掘；多模态模型；人机交互。

\section*{教育背景}
\textbf{澳大利亚国立大学 (ANU)} \hfill 2020.09 -- 2024.09\\
计算机科学博士 (Ph.D.)，导师：邢振昌

\textbf{南京航空航天大学 (NUAA)} \hfill 2017.09 -- 2020.06\\
计算机科学硕士 (M.E.)

\textbf{南京航空航天大学 (NUAA)} \hfill 2013.09 -- 2017.06\\
计算机科学学士 (B.E.)

\section*{发表论文}
\begin{itemize}
  \item \href{https://conf.researchr.org/home/fse-2025}{\textbf{Yanqi Su}, Zhenchang Xing, Chong Wang, Chunyang Chen, Xiwei Xu, Qinghua Lu, Liming Zhu: Automated Soap Opera Testing Directed by LLMs and Scenario Knowledge: Feasibility, Challenges, and Road Ahead. The ACM International Conference on the Foundations of Software Engineering (FSE2025). }
\item \href{https://conf.researchr.org/details/icse-2024/icse-2024-research-track/157/Enhancing-Exploratory-Testing-by-Large-Language-Model-and-Knowledge-Graph}{\textbf{Yanqi Su}, Dianshu Liao, Zhenchang Xing, Qing Huang, Mulong Xie, Qinghua Lu, Xiwei Xu:
Enhancing Exploratory Testing by Large Language Model and Knowledge Graph. The 46th International Conference on Software Engineering (ICSE'24). }
\item \href{https://ieeexplore.ieee.org/document/10172537}{\textbf{Yanqi Su}, Zheming Han, Zhenchang Xing, Xiwei Xu, Liming Zhu, Qinghua Lu:
SoapOperaTG: A Tool for System Knowledge Graph Based Soap Opera Test Generation. The 45th International Conference on Software Engineering (ICSE'23 Tool Demo).}
\item \href{https://dl.acm.org/doi/abs/10.1145/3551349.3556967}{\textbf{Yanqi Su}, Zheming Han, Zhenchang Xing, Xin Xia, Xiwei Xu, Liming Zhu, Qinghua Lu:
Constructing a System Knowledge Graph of User Tasks and Failures from Bug Reports to Support Soap Opera Testing. 37th IEEE/ACM International Conference on Automated Software Engineering (ASE2022).}
\item \href{https://ieeexplore.ieee.org/document/9678574}{\textbf{Yanqi Su}, Zhenchang Xing, Xin Peng, Xin Xia, Chong Wang, Xiwei Xu, Liming Zhu:
Reducing Bug Triaging Confusion by Learning from Mistakes with a Bug Tossing Knowledge Graph. 36th IEEE/ACM International Conference on Automated Software Engineering (ASE2021). \newline
 \hspace{0.5em}\textbf{Distinguished Paper Award}}
\item \href{https://ieeexplore.ieee.org/document/10174001}{\textbf{Yanqi Su}, Zheming Han, Zhipeng Gao, Zhenchang Xing, Qinghua Lu, Xiwei Xu:
Still Confusing for Bug-Component Triaging? Deep Feature Learning and Ensemble Setting to Rescue. The 31st IEEE/ACM International Conference on Program Comprehension (ICPC2023).} 
\item \href{https://ieeexplore.ieee.org/document/8961125}{Yu Zhou, \textbf{Yanqi Su}, Taolue Chen, Zhiqiu Huang, Harald C. Gall, Sebastiano Panichella:
User Review-Based Change File Localization for Mobile Applications. IEEE Trans. on Software Engineering. (TSE)}
\item Zhipeng Gao, \textbf{Yanqi Su}, Xing Hu, Xin Xia: Automating TODO-dropped Methods Detection and Patching. ACM Transactions on Software Engineering and Methodology. (TOSEM)
\end{itemize}

\section*{教学与工作经历}
\begin{itemize}
  \item 博士后研究员，慕尼黑工业大学, Software Engineering \& AI，2024.10 -- 至今（导师：陈春阳）
  \item 讲师，慕尼黑工业大学, 探索性软件测试 (INHN0015, INHN4016, IN0014, IN45111)，2026 夏学期
  \item 讲师，慕尼黑工业大学, 探索性软件测试研讨课，2025 夏学期
  \item 客座讲师，慕尼黑工业大学, 高级软件工程专题，2024 冬学期
  \item 助教，ANU, 软件质量管理, COMP4130，2024 春学期（孙萧育）
  \item 研究助理，复旦大学 CodeWisdom 团队，2020.09 -- 2021.06（彭鑫）
  \item 助教，南京航空航天大学, 软件度量，2017 秋
\end{itemize}

\section*{指导学生}
    \begin{description}
      \item[2025 \textbf{Lazar Minic}, TUM] Master thesis: Bridging Policy and Practice: Leveraging LLMs for Privacy-Aware Exploratory Software Testing
      \item[2024 \textbf{Alper Temel}, TUM] Master thesis: Non-intrusive GUI Element Location Using LLMs
      \item[2024 \textbf{Linsen Mao}, TUM] Bachelor thesis: Soap Opera Testing on Enterprise Resource Planning Software
      \item[2024 \textbf{Bryan Alvin}, TUM] Bachelor thesis: Automated Exploratory Testing of Business Logic in ERP Systems
      \item[2022 \textbf{Zheming Han}, ANU] Master thesis: Reducing Bug Component Triaging Mistakes by Deep Learning
      \item[2017 \textbf{Jian Liu}, NUAA] Bachelor thesis: Code Reviewer Recommendation Based on File Path Analysis
    \end{description}

\section*{资助与奖励}
\begin{itemize}
  \item ANU 校长旅行资助
  \item ACM SIGSOFT CAPS 奖 (FSE/ISSTA 2025)
  \item ACM ICSE 2024 旅行资助
  \item ASE 2021 SIGSOFT 杰出论文奖
  \item ANU 博士学费减免 (2020.09--2024.09)
  \item ANU 大学研究奖学金 (2020.09--2024.09)
  \item NUAA 一等奖学金 (2017.09--2020.04)
  \item NUAA 优秀科研创新奖 (2019.09--2020.04)
  \item NUAA 社会活动贡献奖 (2019.09--2020.04)
  \item Microsoft Research Asia Club Star, 2014
  \item 第2届青年奥运会优秀青年志愿者 (2014.08--2014.09)
\end{itemize}

\section*{学术服务}
期刊审稿人：TOSEM, EMSE；会议程序委员会：MSR2026, ASE2025, ISSRE2025, APSEC2025；会议志愿者：FSE2025.

% \section*{专题演讲}
% \begin{itemize}
%   \item 塑造探索性软件测试的未来：知识图谱与大模型 — HUST(2025.7), CUHK Shenzhen(2025.7), Oxford(2025.6), Aalto(2025.4)
%   \item 基于大模型与场景知识的自动化肥皂剧测试 — TUM(2025.2)
%   \item 构建用户任务与失败的系统知识图 — Google Research Australia(2023.12)
% \end{itemize}

% \section*{推荐人}
% \begin{multicols}{2}
% \textbf{邢振昌 教授}\\
% 澳大利亚国立大学 / CSIRO Data61\\
% Zhenchang.Xing@data61.csiro.au

% \columnbreak

% \textbf{陈春阳 教授}\\
% 慕尼黑工业大学\\
% chun-yang.chen@tum.de
% \end{multicols}

\end{document}
